\chapter{Detailed Derivation of the Expanded Welch--Satterthwaite Test for MI}
\label{chap:expanded}

The Expanded Welch test keeps the estimated MI difference, standard error, and test
statistic from Chapter~\ref{chap:baselines}, but changes the distribution used
to judge that statistic from the standard normal distribution to a Student
\(t\) distribution. To choose its degrees of freedom, this chapter calculates
how much the estimated pointwise-MI variance changes between samples. The name
Expanded Welch refers to this expansion of the pointwise-MI variance as the table
probabilities change.

The derivation first expresses the degrees of freedom in terms of the
sampling mean and variance of the estimated variance contributions. To
approximate those moments, we calculate how one observation changes the
variance estimate, express the sampling error as an average of these
contributions, and take the variance of that average.

\section[Choosing degrees of freedom for the Expanded Welch test]{Choosing degrees of
freedom for\\the Expanded Welch test}
\label{sec:component-df-choice}

\subsection{Connecting the statistic to the Student \texorpdfstring{\(t\)}{t} distribution}

To connect the estimated standard error to the reference distribution,
divide both the estimated MI difference and its standard error by the
population sampling standard deviation of \(\Deltahat\). This rewrites
the shared statistic without changing its value:
\begin{equation*}
  T
  =\frac{\Deltahat/\sqrt{\Var(\Deltahat)}}
         {\sqrt{\SEhat^2/\Var(\Deltahat)}}.
\end{equation*}
The squared denominator is the estimated sampling variance divided by
the population sampling variance. This ratio is the variance term that
the Student \(t\) construction below models.

\begin{samepage}
\medskip
The exact construction uses a numerator with a standard normal
distribution and an independent variance term. If \(Z\) has a standard normal distribution
and \(U\), independently of \(Z\), has a chi-squared distribution with \(\nu\)
degrees of freedom, the scaled ratio follows a Student \(t\) distribution:
\begin{equation}
 \frac{Z}{\sqrt{U/\nu}}\sim t_\nu.
 \label{eq:t-reference-construction}
\end{equation}
\par
\end{samepage}
Under equal MI, \(\Deltahat/\sqrt{\Var(\Deltahat)}\) approaches a standard
normal distribution, as described in Chapter~\ref{chap:baselines}, so it
plays the role of \(Z\) in the approximation. The normalised
squared standard error \(\SEhat^2/\Var(\Deltahat)\) plays the role of
\(U/\nu\), which has a scaled chi-squared distribution. For MI, we model
the estimated variance contributions using scaled chi-squared
distributions with matching sampling means and variances. Because moment
matching does not establish independence between the MI numerator and
the variance estimates, we use the Student \(t\) distribution as an
approximate reference.

\subsection{Moment matching for component degrees of freedom}

The squared standard error adds \(\Vhat(P)/n_P\) and \(\Vhat(Q)/n_Q\).
We first match moments for \(\Vhat(P)\), then show why dividing by
\(n_P\) leaves its component degrees of freedom unchanged.
To obtain the component target \(\nu_{P,\mathrm{mom}}\), Satterthwaite
matches the mean and variance of the nonnegative estimate \(\Vhat(P)\)
with those of a variable following a scaled chi-squared distribution:
\begin{equation}
  \E\{\Vhat(P)\}
  \frac{\chi^2_{\nu_{P,\mathrm{mom}}}}{\nu_{P,\mathrm{mom}}}.
  \label{eq:scaled-chi-working}
\end{equation}

A chi-squared distribution with \(\nu_{P,\mathrm{mom}}\) degrees of freedom
has moments:
\begin{equation}
  \E(\chi^2_{\nu_{P,\mathrm{mom}}})=\nu_{P,\mathrm{mom}},
  \qquad
  \Var(\chi^2_{\nu_{P,\mathrm{mom}}})=2\nu_{P,\mathrm{mom}}.
\end{equation}
The scaling in Equation~\eqref{eq:scaled-chi-working} therefore matches the
mean automatically:
\begin{equation}
  \E\left[
  \E\{\Vhat(P)\}\frac{\chi^2_{\nu_{P,\mathrm{mom}}}}{\nu_{P,\mathrm{mom}}}
  \right]
  =\E\{\Vhat(P)\}.
\end{equation}
The variance of the approximating variable is
\begin{align}
  \Var\left[
  \E\{\Vhat(P)\}\frac{\chi^2_{\nu_{P,\mathrm{mom}}}}{\nu_{P,\mathrm{mom}}}
  \right]
  &=
  \frac{[\E\{\Vhat(P)\}]^2}{\nu_{P,\mathrm{mom}}^2}
  \Var(\chi^2_{\nu_{P,\mathrm{mom}}})\\
  &=
  \frac{2[\E\{\Vhat(P)\}]^2}{\nu_{P,\mathrm{mom}}}.
\end{align}
We now choose \(\nu_{P,\mathrm{mom}}\) so that the approximating variable
also has the same variance as \(\Vhat(P)\). Equating its variance,
calculated above, to the sampling variance of \(\Vhat(P)\) requires
\begin{equation*}
  \frac{2[\E\{\Vhat(P)\}]^2}{\nu_{P,\mathrm{mom}}}
  =\Var\{\Vhat(P)\}.
\end{equation*}
This condition determines the degrees of freedom for the approximating
distribution. Solving for \(\nu_{P,\mathrm{mom}}\) gives
\begin{equation}
  \boxed{
  \nu_{P,\mathrm{mom}}
  =
  \frac{2[\E\{\Vhat(P)\}]^2}
       {\Var\{\Vhat(P)\}}
  }.
  \label{eq:component-df-moments}
\end{equation}
The squared standard error contains \(\Vhat(P)/n_P\), so we apply
moment matching to this contribution as well as to \(\Vhat(P)\).
Because the sample size \(n_P\) is fixed, its reciprocal can be taken
outside the expectation and its square outside the variance.
Applying the same moment-matching calculation to this variance contribution
therefore gives
\begin{align*}
  \frac{2[\E\{\Vhat(P)/n_P\}]^2}
       {\Var\{\Vhat(P)/n_P\}}
  &=\frac{2[\E\{\Vhat(P)\}]^2/n_P^2}
          {\Var\{\Vhat(P)\}/n_P^2}\\
  &=\frac{2[\E\{\Vhat(P)\}]^2}
          {\Var\{\Vhat(P)\}}
   =\nu_{P,\mathrm{mom}}.
\end{align*}
The factors \(1/n_P^2\) cancel, so dividing the estimate by \(n_P\)
leaves its moment-matched degrees of freedom unchanged.

The same calculation applies to the contribution from \(Q\).

To evaluate Equation~\eqref{eq:component-df-moments}, we need the mean and
variance of \(\Vhat(P)\) across repeated samples from \(P\). The first-order
expansion derived in Section~\ref{sec:component-df} gives
\begin{equation*}
  \E\{\Vhat(P)\}\approx V(P).
\end{equation*}
At this order, the change in \(\Vhat(P)\) is linear in the empirical
cell-probability errors. Each empirical cell proportion has expectation
equal to its population probability, so these errors and their linear
contribution have mean zero. Any bias therefore arises from terms beyond
this linear approximation.

With the mean approximated by \(V(P)\), the remaining task is to
calculate \(\Var\{\Vhat(P)\}\). Section~\ref{sec:first-order-mi-background}
obtained the sampling variance of MI by writing its estimation error as
an average of observation-level contributions. We use the same strategy
for the pointwise-MI variance estimate: calculate the contribution of one
observation to its error, then find the variance of the average. The next
section derives that contribution, and Section~\ref{sec:component-df}
shows how averaging it approximates the error in \(\Vhat(P)\) across
whole sampled tables.

\section{How pointwise-MI variance changes with the table}

\subsection{Perturbation and variance sensitivity}

To find the contribution of one observation, we calculate how a small
change in the probability table affects
\begin{equation}
  V(P)
  =\Var_P\{\pmi_P(X,Y)\}.
\end{equation}
Writing this variance in second-moment form identifies the quantities
whose changes we must calculate:
\begin{equation}
  V(P)=\E_P\{\pmi_P(X,Y)^2\}-I(P)^2.
  \label{eq:v-second-moment}
\end{equation}
The first term changes when either the probability weights or the
pointwise-MI values change, and the second changes with MI. We will
differentiate these terms along a path representing one observation.

\begin{samepage}
Fix a cell \((x,y)\) and define a distribution that moves a small amount
of probability toward this cell \citep{hampel1974,vanderVaart1998}:
\begin{equation}
  P_\varepsilon
  =(1-\varepsilon)P+\varepsilon\delta_{(x,y)},
  \label{eq:contamination-path}
\end{equation}
where \(\delta_{(x,y)}\) places probability one on cell \((x,y)\). At
\(\varepsilon=0\), the distribution is \(P\). Multiplying \(P\) by
\(1-\varepsilon\) reduces every original cell probability proportionally,
and adding probability \(\varepsilon\) at \((x,y)\) keeps the total at
one. The selected cell gains probability overall, while all other cells
lose probability. The direction \(\delta_{(x,y)}-P\) is the
difference between the table representing one observation and the
population table, so it describes that observation's contribution to
the empirical-table error.
\par
\end{samepage}

The influence function of \(V(P)\) measures its observation-level
sensitivity and is defined by
\begin{equation}
  g_P(x,y)
  =
  \left.
  \frac{\mathrm d V(P_\varepsilon)}{\mathrm d\varepsilon}
  \right|_{0}.
  \label{eq:g-definition}
\end{equation}
This derivative measures the change in \(V\) per unit increase in the
mixture weight \(\varepsilon\) at zero. The empirical table averages the
one-observation tables, and the first derivative is linear in the table
change, so averaging their sensitivities approximates the error in
\(\Vhat(P)\) at first order. Each observation contributes
\(g_P(x,y)/n_P\) to this error.

\begin{samepage}
For independent observations, the variance of the average gives the approximation derived in
Section~\ref{sec:component-df}:
\begin{equation*}
  \Var\{\Vhat(P)\}
  \approx\frac{\Var_P\{g_P(X,Y)\}}{n_P}.
\end{equation*}
To evaluate the right-hand side, we first need an explicit expression
for \(g_P(x,y)\).
\par
\end{samepage}

Differentiating Equation~\eqref{eq:v-second-moment} separates this
expression into the changes in the second moment and squared MI:
\begin{equation}
  g_P(x,y)
  =
  \left.\frac{\mathrm d}{\mathrm d\varepsilon}
  \E_{P_\varepsilon}\{\pmi_{P_\varepsilon}(X,Y)^2\}
  \right|_{0}
  -2I(P)
  \left.\frac{\mathrm d I(P_\varepsilon)}{\mathrm d\varepsilon}
  \right|_{0}.
  \label{eq:g-decomposition}
\end{equation}
Both terms depend on how pointwise MI changes, which in turn depends on
the joint and marginal probabilities. The selected cell \((x,y)\)
remains fixed, while \((i,j)\) indexes any cell whose probability or
pointwise MI is being evaluated. Along the chosen path,
\begin{equation}
  p_{ij}(\varepsilon)
  =(1-\varepsilon)p_{ij}
  +\varepsilon\mathbf 1\{i=x,j=y\}.
\end{equation}
The joint and marginal probability derivatives at the original distribution
are
\begin{align}
  \left.\frac{\mathrm d p_{ij}(\varepsilon)}{\mathrm d\varepsilon}
  \right|_{0}
  &=\mathbf 1\{i=x,j=y\}-p_{ij},
  \label{eq:joint-perturbation-derivative}\\
  \left.\frac{\mathrm d p_{i+}(\varepsilon)}{\mathrm d\varepsilon}
  \right|_{0}
  &=\mathbf 1\{i=x\}-p_{i+},\\
  \left.\frac{\mathrm d p_{+j}(\varepsilon)}{\mathrm d\varepsilon}
  \right|_{0}
  &=\mathbf 1\{j=y\}-p_{+j}.
\end{align}
The total probability remains one along the path, so summing the
joint-probability derivatives over all cells must give zero:
\begin{equation}
  \sum_{i,j}
  \left.\frac{\mathrm d p_{ij}(\varepsilon)}{\mathrm d\varepsilon}
  \right|_{0}
  =
  \sum_{i,j}
  \{\mathbf 1(i=x,j=y)-p_{ij}\}
  =1-1=0.
\end{equation}

\subsection{Changes in pointwise MI and MI}

The joint and marginal derivatives now determine the change in
pointwise MI. Along the perturbation path,
\begin{equation}
  \pmi_{P_\varepsilon}(i,j)
  =\log p_{ij}(\varepsilon)
  -\log p_{i+}(\varepsilon)
  -\log p_{+j}(\varepsilon).
\end{equation}
\begin{samepage}%
For each logarithm, the derivative is the probability derivative divided by
the probability. Substituting the joint and marginal derivatives gives
\begin{align}
  \left.
  \frac{\mathrm d}{\mathrm d\varepsilon}
  \pmi_{P_\varepsilon}(i,j)
  \right|_{0}
  &={}
  \frac{\mathbf 1\{i=x,j=y\}-p_{ij}}{p_{ij}}
  -\frac{\mathbf 1\{i=x\}-p_{i+}}{p_{i+}}
  -\frac{\mathbf 1\{j=y\}-p_{+j}}{p_{+j}}
  \notag\\*
  &={}
  \frac{\mathbf 1\{i=x,j=y\}}{p_{ij}}
  -\frac{\mathbf 1\{i=x\}}{p_{i+}}
  -\frac{\mathbf 1\{j=y\}}{p_{+j}}
  +1.
  \label{eq:pmi-derivative}
\end{align}
The joint term supplies \(-1\), while subtracting the row and column
terms supplies \(+1\) each. Their sum is the final \(+1\).
\par
\end{samepage}

MI along the path is
\begin{equation}
  I(P_\varepsilon)
  =\sum_{i,j}p_{ij}(\varepsilon)
  \pmi_{P_\varepsilon}(i,j).
\end{equation}

The product rule separates its derivative into the change in probability
weights and the change in pointwise-MI values:
\begin{align}
  \left.\frac{\mathrm d I(P_\varepsilon)}{\mathrm d\varepsilon}
  \right|_{0}
  ={}&
  \sum_{i,j}
  \{\mathbf 1(i=x,j=y)-p_{ij}\}\pmi_P(i,j)
  \notag\\
  &+
  \sum_{i,j}p_{ij}
  \left.\frac{\mathrm d\pmi_{P_\varepsilon}(i,j)}
  {\mathrm d\varepsilon}\right|_{0}.
\end{align}
The first sum is \(\pmi_P(x,y)-I(P)\). Evaluating the second sum using
Equation~\eqref{eq:pmi-derivative}, the selected-cell term contributes
\(p_{xy}/p_{xy}=1\); the row term contributes
\(-\sum_j p_{xj}/p_{x+}=-1\); the column term likewise contributes \(-1\);
and the constant contributes \(\sum_{i,j}p_{ij}=1\). These four
contributions cancel. Hence
\begin{equation}
  \boxed{
  \left.\frac{\mathrm d I(P_\varepsilon)}{\mathrm d\varepsilon}
  \right|_{0}
  =\pmi_P(x,y)-I(P)
  }.
  \label{eq:mi-derivative}
\end{equation}
This is the centred pointwise-MI influence function from
Section~\ref{sec:first-order-mi-background}.

\subsection{Derivative of the pointwise-MI variance}

The MI derivative supplies the second part of
Equation~\eqref{eq:g-decomposition}. To obtain the first part, it remains
to differentiate the second moment,
\begin{equation}
  \E_{P_\varepsilon}\{\pmi_{P_\varepsilon}(X,Y)^2\}
  =\sum_{i,j}p_{ij}(\varepsilon)
  \pmi_{P_\varepsilon}(i,j)^2.
\end{equation}
Applying the product rule to the probability weights and squared
pointwise-MI values gives
\begin{align}
  \left.\frac{\mathrm d}{\mathrm d\varepsilon}
  \E_{P_\varepsilon}\{\pmi_{P_\varepsilon}(X,Y)^2\}
  \right|_{0}
  ={}&
  \sum_{i,j}
  \left.\frac{\mathrm d p_{ij}(\varepsilon)}{\mathrm d\varepsilon}
  \right|_{0}
  \pmi_P(i,j)^2
  \notag\\
  &+
  2\sum_{i,j}p_{ij}\pmi_P(i,j)
  \left.\frac{\mathrm d\pmi_{P_\varepsilon}(i,j)}
  {\mathrm d\varepsilon}\right|_{0}.
  \label{eq:m2-product-rule}
\end{align}
The change in the probability weights contributes
\begin{equation}
  \pmi_P(x,y)^2-\E_P\{\pmi_P(X,Y)^2\}.
  \label{eq:m2-weight-term}
\end{equation}

To express the margin changes in the second sum, define the
probability-weighted row and column means
\begin{align}
  \E_P\{\pmi_P(X,Y)\mid X=x\}
  &=\frac{\sum_j p_{xj}\pmi_P(x,j)}{p_{x+}},\\
  \E_P\{\pmi_P(X,Y)\mid Y=y\}
  &=\frac{\sum_i p_{iy}\pmi_P(i,y)}{p_{+y}}.
\end{align}
To evaluate the sum multiplied by \(2\) in
Equation~\eqref{eq:m2-product-rule}, substitute the pointwise-MI derivative
from Equation~\eqref{eq:pmi-derivative}. The sum has four terms. The
indicators in the first three select the chosen cell, its row, and its column:
\begin{align*}
  \sum_{i,j}p_{ij}\pmi_P(i,j)
  \frac{\mathbf 1\{i=x,j=y\}}{p_{ij}}
  &=\pmi_P(x,y),\\
  \sum_{i,j}p_{ij}\pmi_P(i,j)
  \frac{\mathbf 1\{i=x\}}{p_{i+}}
  &=\frac{\sum_j p_{xj}\pmi_P(x,j)}{p_{x+}},\\
  \sum_{i,j}p_{ij}\pmi_P(i,j)
  \frac{\mathbf 1\{j=y\}}{p_{+j}}
  &=\frac{\sum_i p_{iy}\pmi_P(i,y)}{p_{+y}},\\
  \sum_{i,j}p_{ij}\pmi_P(i,j)&=I(P).
\end{align*}
\begin{samepage}%
The row indicator fixes \(i=x\) while leaving the sum over columns, so
dividing by \(p_{x+}\) gives the conditional mean within row \(x\). The column
term works in the same way. The final term comes from the constant \(+1\)
in the pointwise-MI derivative and averages over the entire table. Keeping
the minus signs on the row and column terms gives
\begin{align}
  &\sum_{i,j}p_{ij}\pmi_P(i,j)
  \left.\frac{\mathrm d\pmi_{P_\varepsilon}(i,j)}
  {\mathrm d\varepsilon}\right|_{0}
  \notag\\
  &\quad=
  \pmi_P(x,y)
  -\frac{\sum_j p_{xj}\pmi_P(x,j)}{p_{x+}}
  -\frac{\sum_i p_{iy}\pmi_P(i,y)}{p_{+y}}
  +I(P)
  \notag\\
  &\quad=
  \pmi_P(x,y)
  -\E_P\{\pmi_P(X,Y)\mid X=x\}
  -\E_P\{\pmi_P(X,Y)\mid Y=y\}
  +I(P).
  \label{eq:m2-score-term}
\end{align}
\par
\end{samepage}

Subtracting the derivative of squared MI from the second-moment derivative
gives \(g_P(x,y)\), as in Equation~\eqref{eq:g-decomposition}.
First combine the probability-weight change in
Equation~\eqref{eq:m2-weight-term} with the derivative of \(-I(P)^2\):
\begin{align*}
 &\pmi_P(x,y)^2-\E_P\{\pmi_P(X,Y)^2\}
       -2I(P)\{\pmi_P(x,y)-I(P)\}\\*
 &\quad=\pmi_P(x,y)^2-\{V(P)+I(P)^2\}
       -2I(P)\pmi_P(x,y)+2I(P)^2\\*
 &\quad=\pmi_P(x,y)^2-2I(P)\pmi_P(x,y)+I(P)^2-V(P)\\*
 &\quad=\{\pmi_P(x,y)-I(P)\}^2-V(P).
\end{align*}
\begin{samepage}%
The second line uses
\(\E_P\{\pmi_P(X,Y)^2\}=V(P)+I(P)^2\); the last line completes the
square. Adding twice the expression in Equation~\eqref{eq:m2-score-term}
gives
\begin{equation}
\boxed{
\begin{aligned}
g_P(x,y)
={}&\{\pmi_P(x,y)-I(P)\}^2-V(P)\\
&+2\bigl[
\pmi_P(x,y)
-\E_P\{\pmi_P(X,Y)\mid X=x\}\\
&\hspace{18mm}
-\E_P\{\pmi_P(X,Y)\mid Y=y\}
+I(P)
\bigr].
\end{aligned}
}
\label{eq:g-final}
\end{equation}
The first line describes the change with the pointwise-MI values held
fixed, while the bracket accounts for changes in those values as the cell
probabilities and their row and column totals change. The conditional
means collect the effects of the changing totals, so both parts are needed
to describe how the estimated variance varies across samples.
\par
\end{samepage}

\section{Component degrees of freedom and final test}
\label{sec:component-df}

\subsection{Sampling error of the variance estimate}

We now express \(\Vhat(P)-V(P)\) as an average of the sensitivities
derived in the preceding section. The sampling mean and variance of
this average will supply the two inputs to the component degrees of
freedom.

Let \(Z_a^{(P)}\) denote observation \(a\) from population \(P\). It can
be represented by a table that places probability one in its cell,
denoted \(\delta_{Z_a^{(P)}}\). Relative to
\(P\), that observation contributes the table change
\(\delta_{Z_a^{(P)}}-P\). The empirical table averages these one-observation
tables, so its difference from \(P\) is
\begin{equation}
  \widehat P-P
  =\frac{1}{n_P}\sum_{a=1}^{n_P}
  \left\{\delta_{Z_a^{(P)}}-P\right\}.
\end{equation}
The plug-in estimate \(\Vhat(P)\) evaluates \(V\) at the empirical table,
so \(\Vhat(P)=V(\widehat P)\). Near a positive population table, \(V\)
is smooth, so we approximate its change by a first-order Taylor
expansion. For a differentiable scalar function
expanded at zero, the general formula is
\begin{equation*}
  f(t)\approx f(0)+t f'(0).
\end{equation*}
Here set \(f(t)=V\{P+t(\widehat P-P)\}\). As the scalar \(t\) increases,
the table moves along the straight line from the population table \(P\) at
\(t=0\) to the empirical table \(\widehat P\) at \(t=1\). Thus \(V\) is expanded at \(P\), the
new table is \(\widehat P\), and the table difference is
\(\widehat P-P\). The size of this change is determined by the
cell-probability errors \(\widehat P-P\), while \(t\) specifies a
position along the path. The approximation requires these probability
errors to be small. Since \(f(0)=V(P)\) and \(f(1)=\Vhat(P)\), evaluating
the Taylor formula at \(t=1\) and subtracting \(f(0)\) gives
\begin{equation*}
  \Vhat(P)-V(P)
  \approx f'(0)
  =\left.\frac{\mathrm d}{\mathrm dt}
    V\{P+t(\widehat P-P)\}\right|_{t=0}.
\end{equation*}
The omitted terms are quadratic or higher order in the cell-probability
errors, so the approximation can be useful at \(t=1\) when these errors
are small.

At fixed \(P\), the linear Taylor term is a weighted sum of the
cell-probability changes, with coefficients evaluated at \(P\). These
coefficients stay fixed when the direction of perturbation changes, so
the derivative of an average table change is the average of the
individual derivatives \citep[Chapter~20]{vanderVaart1998}. Substituting
the empirical-table identity above therefore gives
\begin{equation*}
  \left.\frac{\mathrm d}{\mathrm dt}
    V\{P+t(\widehat P-P)\}\right|_{t=0}
  =\frac{1}{n_P}\sum_{a=1}^{n_P}
    \left.\frac{\mathrm d}{\mathrm dt}
    V\{P+t(\delta_{Z_a^{(P)}}-P)\}\right|_{t=0}.
\end{equation*}
Each term follows the one-observation perturbation
\(P+t(\delta_{Z_a^{(P)}}-P)=(1-t)P+t\delta_{Z_a^{(P)}}\), so its
derivative is \(g_P(Z_a^{(P)})\) by Equation~\eqref{eq:g-definition}.
Combining this derivative identity with the Taylor approximation gives
\begin{equation}
  \Vhat(P)-V(P)
  \approx
  \frac{1}{n_P}
  \sum_{a=1}^{n_P}g_P(Z_a^{(P)}).
  \label{eq:vhat-linearization}
\end{equation}
The variance-estimation error now has the same sample-average form as the
MI error in Equation~\eqref{eq:mi-error-as-average-background}, with
\(g_P\) supplying the observation-level contribution.

\subsection{Sampling mean and variance of the estimate}

The linearised variance-estimation error averages the sensitivities
\(g_P\), so its mean is zero if their mean is zero. We verify this
property by averaging the two parts of Equation~\eqref{eq:g-final}.
The first part has mean
\begin{equation}
  \E_P\left[\{\pmi_P(X,Y)-I(P)\}^2-V(P)\right]
  =V(P)-V(P)=0.
\end{equation}
For the bracket, the row and column conditional means are averaged by
weighting each row or column by its probability:
\begin{align}
  \E_P\left[\E_P\{\pmi_P(X,Y)\mid X\}\right]
  &=\sum_x p_{x+}
    \frac{\sum_y p_{xy}\pmi_P(x,y)}{p_{x+}}
    \notag\\
  &=\sum_{x,y}p_{xy}\pmi_P(x,y)=I(P),\\
  \E_P\left[\E_P\{\pmi_P(X,Y)\mid Y\}\right]
  &=\sum_y p_{+y}
    \frac{\sum_x p_{xy}\pmi_P(x,y)}{p_{+y}}
    \notag\\
  &=\sum_{x,y}p_{xy}\pmi_P(x,y)=I(P).
\end{align}
Each marginal probability cancels the denominator of its conditional mean,
leaving the probability-weighted sum over all cells. The bracket therefore
also has mean zero, and the full sensitivity satisfies
\begin{equation}
  \E_P\{g_P(X,Y)\}
  =0+2\{I(P)-I(P)-I(P)+I(P)\}=0.
\end{equation}
\begin{samepage}%
Taking expectations in Equation~\eqref{eq:vhat-linearization} and using
the zero mean of \(g_P\) gives
\begin{align*}
  \E\{\Vhat(P)\}-V(P)
  &\approx\frac{1}{n_P}\sum_{a=1}^{n_P}
    \E_P\{g_P(Z_a^{(P)})\}\\
  &=0.
\end{align*}
\par
\end{samepage}
This establishes the first-order mean approximation used in
Section~\ref{sec:component-df-choice}. The other required moment depends
on the variation of \(g_P\) across possible observations. Define this
population variance by
\begin{equation}
  \tau^2(P)=\Var_P\{g_P(X,Y)\}.
  \label{eq:tau-definition}
\end{equation}
Since \(g_P\) has mean zero, \(\tau^2(P)\) measures variation in the
sensitivities around zero. The independent observations give
independent sensitivity contributions, so their covariance terms vanish
in the following first-order sampling-variance approximation:
\begin{align}
  \Var\{\Vhat(P)\}
  &\approx
  \frac{1}{n_P^2}
  \sum_{a=1}^{n_P}
  \Var_P\{g_P(Z_a^{(P)})\}\\*
  &=\frac{\tau^2(P)}{n_P}.
  \label{eq:vhat-sampling-variance}
\end{align}

\subsection{Population and estimated component degrees of freedom}

The mean and variance approximations now supply both inputs to
Equation~\eqref{eq:component-df-moments}. To obtain positive, finite
component degrees of freedom, we require positive sensitivity variances:
\begin{assumption}[Positive variances of observation-level sensitivities]
For the first-order construction of the Expanded Welch test,
\(\tau^2(P)>0\) and \(\tau^2(Q)>0\).
\end{assumption}
Substituting the first-order mean and variance approximations into
Equation~\eqref{eq:component-df-moments} gives
\begin{equation*}
  \nu_{P,\mathrm{mom}}
  \approx\frac{2V(P)^2}{\tau^2(P)/n_P}
  =\frac{2n_PV(P)^2}{\tau^2(P)}.
\end{equation*}
We define this population first-order approximation by
\begin{equation}
  \boxed{
  \nu_{P,\mathrm{FO}}
  =
  \frac{2n_PV(P)^2}{\tau^2(P)}
  }.
  \label{eq:population-component-df}
\end{equation}
To interpret Equation~\eqref{eq:population-component-df}, divide the
sampling variance of \(\Vhat(P)\) by \(V(P)^2\), the square of the
population pointwise-MI variance:
\begin{equation}
  \frac{\Var\{\Vhat(P)\}}{V(P)^2}
  \approx \frac{\tau^2(P)}{n_PV(P)^2},
\end{equation}
so the component degrees of freedom are approximately twice the reciprocal
of this relative variance. When \(\Vhat(P)\) fluctuates little relative to
\(V(P)\), the degrees of freedom are large and the Student \(t\)
distribution is close to the standard normal distribution. Greater relative
fluctuation instead gives fewer degrees of freedom and a larger rejection
cutoff.
Positive \(V(P)\) does not guarantee positive \(\tau^2(P)\): table changes
can have no first-order effect on \(V(P)\). In that case, describing the
sampling variation of \(\Vhat(P)\) requires a higher-order approximation.

\medskip
\thesisminiheading{Table-based component degrees of freedom.}

The population inputs \(V(P)\) and \(\tau^2(P)\) are unknown, so the
test estimates them from the observed table. We replace every population
quantity in Equation~\eqref{eq:g-final} by its empirical counterpart and sum over
occupied cells, since pointwise MI is undefined at an empirical zero.
For each occupied cell,
\begin{equation}
\begin{aligned}
\widehat g_P(i,j)
={}&\{\widehat\pmi_P(i,j)-\Ihat(P)\}^2-\Vhat(P)\\
&+2\left[
\widehat\pmi_P(i,j)
-\frac{\sum_{j'}\widehat p_{ij'}\widehat\pmi_P(i,j')}
       {\widehat p_{i+}}
-\frac{\sum_{i'}\widehat p_{i'j}\widehat\pmi_P(i',j)}
       {\widehat p_{+j}}
+\Ihat(P)
\right].
\end{aligned}
\label{eq:g-hat}
\end{equation}
An empty cell has zero weight in these sums, so no logarithm or sensitivity
needs to be evaluated there. Although the empirical sensitivities have
probability-weighted mean zero under exact substitution, their variance
is centred explicitly:
\begin{equation}
  \widehat\tau^2(P)
  =\sum_{i,j}\widehat p_{ij}
  \left\{\widehat g_P(i,j)
  -\sum_{k,l}\widehat p_{kl}\widehat g_P(k,l)\right\}^2.
  \label{eq:tau-hat}
\end{equation}
The weighted variance \(\widehat\tau^2(P)\) estimates the population
variance of the observation-level sensitivities. By
Equation~\eqref{eq:vhat-sampling-variance}, \(\widehat\tau^2(P)/n_P\)
then estimates the sampling variance of \(\Vhat(P)\). Replacing \(V(P)\) and
\(\tau^2(P)\) by their estimates in
Equation~\eqref{eq:population-component-df} gives the estimated component
degrees of freedom:
\begin{equation}
  \boxed{
  \widehat\nu_P
  =\frac{2n_P\Vhat(P)^2}{\widehat\tau^2(P)}
  }.
  \label{eq:nu-p-hat}
\end{equation}
Repeating the same calculation for the second table gives
\(\widehat\nu_Q\).

\medskip
Table~\ref{tab:component-df-roles} distinguishes the sampling-moment target,
its population first-order approximation, and the sample-based estimate.
The target and population approximation are fixed for a population and
sample size, whereas \(\widehat\nu_P\) changes with each sampled table.

\begin{table}[htbp]
\centering
\caption{Component degrees of freedom and their inputs.}
\label{tab:component-df-roles}
\thesistableformat
\begin{tabularx}{\textwidth}{@{}l>{\raggedright\arraybackslash}X>{\raggedright\arraybackslash}X@{}}
\toprule
\textbf{Quantity} & \textbf{Role} & \textbf{Inputs}\\
\midrule
\(\nu_{P,\mathrm{mom}}\) & Sampling-moment target &
\(\E\{\Vhat(P)\}\) and \(\Var\{\Vhat(P)\}\) at fixed \(P,n_P\)\\
\addlinespace
\(\nu_{P,\mathrm{FO}}\) & Population first-order approximation to the target &
Population \(V(P)\), \(\tau^2(P)\), and \(n_P\)\\
\addlinespace
\(\widehat\nu_P\) & Component estimate used by Expanded Welch &
\(\Vhat(P)\), \(\widehat\tau^2(P)\), and \(n_P\) from one sample\\
\bottomrule
\end{tabularx}
\end{table}

\subsection{Combined degrees of freedom and p-value}

The two component degrees of freedom must now be combined to choose one
reference distribution for \(T\). For component inputs \(\nu_P\) and
\(\nu_Q\), the Welch--Satterthwaite equation is
\citep{satterthwaite1946,welch1947}:
\begin{equation}
  \nu
  =
  \frac{
  \left\{\Vhat(P)/n_P+\Vhat(Q)/n_Q\right\}^2
  }{
  \left\{\Vhat(P)/n_P\right\}^2/\nu_P
  +\left\{\Vhat(Q)/n_Q\right\}^2/\nu_Q
  }.
  \label{eq:ws-general}
\end{equation}
To derive this rule, we apply moment matching to the full estimated
variance,
\begin{equation*}
  \SEhat^2=\frac{\Vhat(P)}{n_P}+\frac{\Vhat(Q)}{n_Q}.
\end{equation*}
Linearity of expectation gives its sampling mean:
\begin{equation*}
  \E(\SEhat^2)
  =\frac{\E\{\Vhat(P)\}}{n_P}
   +\frac{\E\{\Vhat(Q)\}}{n_Q}.
\end{equation*}
The two pointwise-MI variance estimates are computed from independent
samples, so their covariance is zero and their contributions to
\(\Var(\SEhat^2)\) add. Rearranging
Equation~\eqref{eq:component-df-moments} expresses each component's sampling
variance in terms of its sampling mean and moment-matched degrees of
freedom. Applying this identity to both components gives
\begin{align*}
  \Var(\SEhat^2)
  &=\frac{\Var\{\Vhat(P)\}}{n_P^2}
    +\frac{\Var\{\Vhat(Q)\}}{n_Q^2}\\*
  &=2\left\{
    \frac{[\E\{\Vhat(P)\}/n_P]^2}{\nu_{P,\mathrm{mom}}}
    +\frac{[\E\{\Vhat(Q)\}/n_Q]^2}{\nu_{Q,\mathrm{mom}}}
    \right\}.
\end{align*}
The combined moment-matched degrees of freedom are twice the squared mean
of \(\SEhat^2\) divided by its variance. Substituting the expressions above
and cancelling the common factor \(2\) gives
\begin{equation*}
  \frac{2[\E(\SEhat^2)]^2}{\Var(\SEhat^2)}
  =\frac{
    \left\{\E\{\Vhat(P)\}/n_P+\E\{\Vhat(Q)\}/n_Q\right\}^2
    }{
    [\E\{\Vhat(P)\}/n_P]^2/\nu_{P,\mathrm{mom}}
    +[\E\{\Vhat(Q)\}/n_Q]^2/\nu_{Q,\mathrm{mom}}
    }.
\end{equation*}
This establishes the combining rule in Equation~\eqref{eq:ws-general} for
population moment inputs. For the observed tables, we replace the sampling
means \(\E\{\Vhat(P)\}\) and \(\E\{\Vhat(Q)\}\) by \(\Vhat(P)\) and
\(\Vhat(Q)\), and the component targets \(\nu_{P,\mathrm{mom}}\) and
\(\nu_{Q,\mathrm{mom}}\) by \(\widehat\nu_P\) and \(\widehat\nu_Q\).
These substitutions give the estimated combined degrees of freedom:
\begin{equation}
  \boxed{
  \widehat\nu_{\mathrm{expanded}}
  =
  \frac{
  \left\{\Vhat(P)/n_P+\Vhat(Q)/n_Q\right\}^2
  }{
  \left\{\Vhat(P)/n_P\right\}^2/\widehat\nu_P
  +
  \left\{\Vhat(Q)/n_Q\right\}^2/\widehat\nu_Q
  }
  }.
  \label{eq:expanded-df}
\end{equation}
The two-sided p-value for the Expanded Welch test is
\begin{equation}
  p_{\mathrm{expanded}}
  =2\left\{1-F_t(|T|;\widehat\nu_{\mathrm{expanded}})\right\},
  \label{eq:expanded-p}
\end{equation}
where \(F_t(x;\nu)\) denotes the cumulative distribution function of a
Student \(t\) distribution with \(\nu\) degrees of freedom, evaluated at \(x\).

For a fixed finite \(T\), the two-sided p-value from the Student \(t\)
distribution is at least as large as the Wald p-value from the standard
normal distribution.
We call a test calculation successful when all required quantities satisfy
the numerical conditions in Appendix~\ref{sec:validity-rules}. Every
successful calculation for the Expanded Welch test also satisfies the Wald test's
calculation requirements. At the same p-value threshold, every rejection by the
Expanded Welch test is therefore also a rejection by the Wald test, including when unsuccessful calculations
count as non-rejections in the reported rates. At the stated significance
level, the Expanded Welch test cannot raise rejection where the Wald test is already
conservative.

\section{Worked example and implementation}

\subsection{Worked comparison}

Consider the two observed \(3\times3\) tables
\begin{equation}
  N^P=
  \begin{pmatrix}
  35&10&5\\
  10&20&5\\
  5&5&10
  \end{pmatrix},
  \qquad
  N^Q=
  \begin{pmatrix}
  30&5&5\\
  5&30&5\\
  5&5&10
  \end{pmatrix}.
  \label{eq:worked-tables}
\end{equation}
Their totals are \(n_P=105\) and \(n_Q=100\). In \(N^P\), cell
\((1,1)\) has count 35, and its row and column each total 50. Its empirical
pointwise MI is therefore
\begin{equation*}
 \widehat\pmi_P(1,1)
 =\log\left\{\frac{35/105}{(50/105)(50/105)}\right\}
 =\log(1.47)\simeq0.38526.
\end{equation*}
This cell contributes \((35/105)(0.38526)\) to the plug-in MI estimate.
Summing the contributions from all nine cells gives
\(\Ihat(P)=0.13693\); the same calculation for \(N^Q\) gives
\(\Ihat(Q)=0.25848\) nats. Since \((3-1)(3-1)=4\), the corrected difference is
\begin{equation}
  \Deltahat
  =\left(0.13693-\frac{4}{2(105)}\right)
  -\left(0.25848-\frac{4}{2(100)}\right)
  =-0.12059.
\end{equation}

The variance calculation weights each squared deviation from the estimated
MI by its cell probability. Cell \((1,1)\), for example, contributes
\((35/105)(0.38526-0.13693)^2\simeq0.02056\) to \(\Vhat(P)\).
Summing over all cells gives \(\Vhat(P)=0.27189\) and
\(\Vhat(Q)=0.43371\). Hence
\begin{equation}
  \SEhat
  =\left(\frac{0.27189}{105}+\frac{0.43371}{100}\right)^{1/2}
  =0.08323,
  \qquad
  T=-1.44898.
\end{equation}
The Wald test gives \(p=0.14734\).

For the Expanded Welch test, the same cell illustrates the sensitivity calculation.
The first row's pointwise-MI values are 0.38526, \(-0.51083\), and
\(-0.64436\), weighted by its counts 35, 10, and 5. Its conditional mean is
\begin{equation*}
 \frac{35(0.38526)+10(-0.51083)+5(-0.64436)}{50}
 \simeq0.10308.
\end{equation*}
Because this count table is symmetric, the first column has the same
conditional mean. Substituting these values into Equation~\eqref{eq:g-hat} gives
\begin{align*}
 \widehat g_P(1,1)
 &\simeq (0.38526-0.13693)^2-0.27189\\*
 &\quad+2\{0.38526-0.10308-0.10308+0.13693\}\\*
 &\simeq0.4218.
\end{align*}
The probability-weighted mean of the sensitivities is zero, so this cell
contributes \((35/105)(0.4218)^2\simeq0.0593\) to
\(\widehat\tau^2(P)\). Repeating the calculation for all cells in both
tables gives \(\widehat\tau^2(P)=0.93212\) and
\(\widehat\tau^2(Q)=0.82035\). Substitution into
Equation~\eqref{eq:nu-p-hat} gives component degrees of freedom
\begin{equation}
  \widehat\nu_P=16.65,
  \qquad
  \widehat\nu_Q=45.86.
\end{equation}
Combining them with Equation~\eqref{eq:expanded-df} gives
\(\widehat\nu_{\mathrm{expanded}}=59.03\) and \(p=0.15264\).
Neither method rejects at the 5\% level.

\pagebreak[3]
\subsection{Operational algorithm and complexity}
\label{sec:expanded-complexity}

Table~\ref{tab:expanded-algorithm} summarises one comparison.

\begin{table}[htbp]
\centering
\caption{Operational Expanded Welch--Satterthwaite algorithm.}
\label{tab:expanded-algorithm}
\thesistableformat
\begin{tabularx}{\textwidth}{@{}c>{\raggedright\arraybackslash}X@{}}
\toprule
\textbf{Step} & \textbf{Calculation} \\
\midrule
1 & Validate the two aligned count tables and their specified numbers of row and column categories; independence of samples is a design assumption. \\
\addlinespace[0.4em]
2 & Compute empirical probabilities, pointwise MI, plug-in MI, and the leading bias-corrected difference. \\
\addlinespace[0.4em]
3 & Compute \(\Vhat(P)\), \(\Vhat(Q)\), the shared standard error, and \(T\). \\
\addlinespace[0.4em]
4 & Compute the row and column pointwise-MI means needed in \(\widehat g_P\) and \(\widehat g_Q\). \\
\addlinespace[0.4em]
5 & Compute \(\widehat\tau^2(P)\), \(\widehat\tau^2(Q)\), and the two component degrees of freedom. \\
\addlinespace[0.4em]
6 & Combine the components using Equation~\eqref{eq:expanded-df}. \\
\addlinespace[0.4em]
7 & Return the p-value from the Student \(t\) distribution, effective degrees of freedom, and results of the calculation checks. \\
\bottomrule
\end{tabularx}
\end{table}

Each empirical probability, pointwise-MI value, row mean, column mean, and
variance contribution is calculated by a fixed number of table scans. The
time complexity is therefore \(O(rc)\), and the working memory is also
\(O(rc)\).

\section{Behaviour and limitations}
\label{sec:near-independence-mechanism}

\subsection{Near independence}

Consider a population table that approaches independence while its row
and column margins stay fixed. At independence, every pointwise-MI value
is zero. Since the linear MI error in
Equation~\eqref{eq:mi-error-from-cell-errors} weights each cell-probability
error by its population pointwise MI, this linear contribution vanishes.
Sampling errors can still create apparent association through products of
cell-probability errors, which are retained by a second-order expansion.
Near independence, the small linear weights allow these quadratic terms
to contribute substantially to the MI estimation error.

The fixed-margin expansion in Appendix~\ref{sec:fixed-margin-local-derivation}
relates MI, the pointwise-MI variance \(V(P)\), and the variance of the
sensitivities \(\tau^2(P)\) under the assumptions stated there. Close to
independence along that path, with fixed table dimensions and positive
cell probabilities, the expansion gives
\begin{equation}
 V(P)\approx 2I(P),
 \qquad \tau^2(P)\approx4V(P).
 \label{eq:near-independent-moments}
\end{equation}
Substituting these local approximations into
the component degrees-of-freedom formula gives
\begin{equation}
 \nu_{P,\mathrm{FO}}=\frac{2n_PV(P)^2}{\tau^2(P)}
 \approx\frac{n_PV(P)}{2}
 \approx n_PI(P).
 \label{eq:near-independent-df}
\end{equation}
At a fixed sample size, this population first-order value tends to zero
as positive population MI approaches zero. When \(\Vhat(P)\) and
\(\widehat\tau^2(P)\) estimate their population counterparts accurately
enough for the component formula to be reliable, \(\widehat\nu_P\) follows this
prediction: smaller MI gives fewer degrees of freedom and a larger cutoff
from the Student \(t\) distribution.

The empirical table, however, need not be independent even when the
population is independent. Random differences between its observed cell
proportions and the products of its observed margins produce nonzero
pointwise-MI values in occupied cells. Both \(\Vhat(P)\) and \(\widehat\tau^2(P)\) are
calculated from these values, so sampling-induced association enters the
estimated degrees of freedom. At a fixed sample size, these estimated
inputs need not track the shrinking population quantities accurately
enough for \(\widehat\nu_P\) to follow \(n_PI(P)\). Consequently, the
estimated degrees of freedom need not approach zero with population MI.
Appendix~\ref{sec:independence-boundary} retains the quadratic sampling
terms and derives their contribution to the estimated degrees of freedom.

\begin{samepage}
\subsection{Theoretical status and calculation requirements}

Equations~\eqref{eq:mi-derivative} and \eqref{eq:g-final} are exact
derivatives for positive population tables. Using these derivatives to
describe finite sampling errors gives first-order approximations for the
MI estimate and the estimated pointwise-MI variance. Moment matching then
approximates each variance contribution by a variable with a scaled
chi-squared distribution, as in Equation~\eqref{eq:scaled-chi-working}.
The exact construction for a Student \(t\) distribution in
Section~\ref{sec:component-df-choice} also requires a numerator with a
standard normal distribution and a squared denominator with a scaled
chi-squared distribution. The numerator and squared denominator must be
independent. Here, the corrected MI
difference and its estimated standard error come from the same tables, so
an observation can change both quantities. Matching the variance moments
does not establish the required
distributional or independence conditions, and the Student \(t\) distribution
for \(T\) therefore remains a finite-sample approximation.

For fixed population tables satisfying the stated assumptions, including
positive \(V\) and \(\tau^2\), the estimated degrees of freedom grow as
the sample sizes tend to infinity, bringing the Student \(t\) distribution
closer to the standard normal distribution used by the Wald test. This large-sample behaviour holds the
populations fixed, whereas the near-independence comparison above reduces
population MI at a fixed sample size.
\par\nopagebreak[4]

Both tests need a finite corrected MI difference and a positive, finite
estimated standard error, while the Expanded Welch test also needs positive, finite component
and combined degrees of freedom. If a requirement is not met, the calculation
is unsuccessful and cannot be used to draw a test conclusion.
The experiment applies the numerical calculation requirements in
Appendix~\ref{sec:validity-rules} and assesses whether false-positive rates
are close to 5\%.
\par
\end{samepage}
